\section{Related Work}
\label{sec:related}

\subsection{Data-Driven Time-Series Forecasting}
Modern time-series forecasting increasingly relies on deep sequence models in addition to classical statistical and tree-based methods. Transformer-based forecasting has been strengthened by patching and channel-independent processing, as exemplified by PatchTST \cite{nie2022patchtst}. State-space sequence models provide another route to long-context modeling: Mamba introduces input-dependent selective state-space dynamics with linear sequence-length scaling \cite{gu2023mamba}. These model families are why the frozen comparison set in this study includes attention-based and convolutional temporal baselines, although the specific published architectures cited here are not themselves among them. Predictive skill alone, in any case, does not establish whether a forecast should be trusted under non-stationarity.

\subsection{Kolmogorov--Arnold Networks for Temporal Prediction}
Kolmogorov--Arnold Networks (KANs) replace the fixed scalar nonlinearities of conventional multilayer perceptrons with learnable univariate functions on edges, yielding a functional representation that can be visualized and inspected \cite{liu2025kan}. Early work applying KANs to time-series analysis reported competitive forecasting performance with compact parameterizations \cite{vacarubio2024kan}. These developments motivate KANs as an interpretable-by-design forecasting backbone. The present study asks a distinct question: whether a temporal KAN representation can also provide useful evidence about forecast reliability under temporal and distributional shift. Accordingly, interpretability is evaluated through functional stability rather than assumed from architecture alone.

\subsection{Conformal Uncertainty Quantification for Time Series}
Conformal prediction provides a model-agnostic mechanism for calibrating predictive sets from empirical nonconformity scores. Conformalized quantile regression combines learned conditional quantiles with conformal calibration to improve interval adaptivity while retaining marginal coverage guarantees under the corresponding assumptions \cite{romano2019cqr}. Weighted conformal methods have also been developed for covariate shift \cite{tibshirani2019covariate}. For sequential data, conformal time-series forecasting extends inductive conformal ideas to multi-horizon forecasting \cite{stankeviciute2021conformal}, while Adaptive Conformal Inference updates the effective miscoverage level online to address changing distributions \cite{gibbs2021adaptive}. Subsequent work studied adaptive conformal prediction specifically for dependent time series \cite{zaffran2022adaptive}, and control-inspired conformal updates have further targeted systematic temporal errors and distribution shifts \cite{angelopoulos2023pid}. These methods motivate our comparison of static, rolling, and adaptive calibration.

\subsection{Distribution Shift, Reliability, and Selective Prediction}
Selective prediction allows a model to abstain on low-confidence samples, explicitly trading prediction coverage for reduced risk \cite{geifman2017selective}. SelectiveNet further showed that rejection can be incorporated directly into deep models and evaluated through risk--coverage behaviour \cite{geifman2019selectivenet}. Distribution-shift monitoring provides a complementary capability by detecting deviations from a reference operating regime; recent work has connected shift detection with selective-prediction principles in streaming settings \cite{barshalom2023shift}. In \method{}, interval width and representation shift are treated as candidate evidence sources. Crucially, these signals are not assumed to measure reliability: they must demonstrate association with realized forecast error and improved risk--coverage trade-offs on held-out data, and the fusion weight that combines them is selected on the calibration partition rather than frozen in advance.
